% Illustrative application background, not an experimental deployment.
% This scene explains the observation and communication requirements;
% algorithm relationships are given separately in research-framework.tex.
\begin{tikzpicture}[
  x=1mm,y=1mm,>=Stealth,
  font=\fontsize{9bp}{11.5bp}\selectfont,
  bgnote/.style={align=center,fill=white,inner sep=.6mm},
  bgview/.style={draw=black!35,line width=.4pt,densely dotted},
  bgdata/.style={->,draw=black!85,line width=.8pt},
  bgfeedback/.style={->,draw=black!75,line width=.65pt,
    dash pattern=on 2.5pt off 2pt},
  pics/bgdrone/.style={code={
    \draw[black!85,line width=.7pt]
      (-3.4,-1.6)--(3.4,1.6) (-3.4,1.6)--(3.4,-1.6);
    \foreach \px/\py in {-4.1/-1.9,-4.1/1.9,4.1/-1.9,4.1/1.9}
      \draw[black!80,fill=white,line width=.6pt]
        (\px,\py) ellipse[x radius=1.8mm,y radius=.75mm];
    \draw[black!85,fill=black!18,line width=.6pt]
      (0,0) ellipse[x radius=1.45mm,y radius=1mm];
    \draw[black!80,line width=.5pt] (-.8,-.7)--(-.8,-1.6)
      (.8,-.7)--(.8,-1.6);
    \draw[black!80,fill=white,line width=.45pt]
      (-.55,-1.7) rectangle (.55,-2.5);
  }}
]
\path[use as bounding box] (0,0) rectangle (148,88);
\node[anchor=west,font=\fontsize{9bp}{11.5bp}\selectfont\bfseries]
  at (1,85.5) {灾后道路观测（应用示例）};

% Road and the partially overlapping fields of view, in oblique projection.
\filldraw[fill=black!2,draw=black!35,line width=.45pt]
  (3,27)--(91,27)--(110,46)--(22,46)--cycle;
\filldraw[fill=black!12,draw=black!45,line width=.45pt]
  (8,31)--(96,31)--(104,39)--(16,39)--cycle;
\draw[white,line width=.9pt,dash pattern=on 3mm off 2mm]
  (12,35)--(100,35);
\fill[black,opacity=.025]
  (20,62.5)--(7,29)--(46,29)--(61,44)--(22,44)--cycle;
\fill[black,opacity=.025]
  (66,58.5)--(44,29)--(90,29)--(105,44)--(59,44)--cycle;
\draw[black!45,line width=.45pt,densely dotted]
  (7,29)--(46,29)--(61,44)--(22,44)--cycle
  (44,29)--(90,29)--(105,44)--(59,44)--cycle;
\foreach \p in {(7,29),(46,29),(22,44)}
  \draw[bgview] (20,62.5)--\p;
\foreach \p in {(44,29),(90,29),(105,44)}
  \draw[bgview] (66,58.5)--\p;

% A building occludes part of the road from the left-hand viewpoint.
\filldraw[fill=black!10,draw=black!65,line width=.5pt]
  (42,38)--(52,38)--(52,47)--(42,47)--cycle;
\filldraw[fill=black!23,draw=black!65,line width=.5pt]
  (52,38)--(58,44)--(58,53)--(52,47)--cycle;
\filldraw[fill=black!5,draw=black!65,line width=.5pt]
  (42,47)--(52,47)--(58,53)--(48,53)--cycle;
\node[font=\fontsize{8bp}{10bp}\selectfont] at (47,42.5) {建筑};

% Water on the road is a concrete observation target, with no claimed result.
\filldraw[fill=black!22,draw=black!55,line width=.45pt]
  (67,33.2) .. controls (65,35) and (69,37) .. (73,37)
  .. controls (77,38.8) and (85,37.3) .. (86,35.4)
  .. controls (89,33.4) and (84,32) .. (79,32.5)
  .. controls (74,31.8) and (70,32) .. cycle;
\draw[white,line width=.65pt]
  (70,34.5) .. controls (73,35.5) and (76,33.8) .. (79,34.8)
  (77,36) .. controls (79,36.7) and (82,35.5) .. (84,36);
\node[bgnote,inner sep=.3mm] at (24,28.1) {道路入口};
\node[bgnote,inner sep=.3mm] at (78,28.1) {中段积水};
\node[font=\fontsize{8.5bp}{10.5bp}\selectfont,align=center]
  at (53,22.5) {不同视角提供局部观测，部分区域重复覆盖};

% Observing UAVs and the aerial relay.
\pic at (20,65) {bgdrone};
\pic at (66,61) {bgdrone};
\pic at (92,74) {bgdrone};
\node[bgnote] at (20,56.6) {观测无人机 A};
\node[bgnote] at (76,52.8) {观测无人机 B};
\node[bgnote] at (92,82) {中继无人机};

% Data travels toward the receiver; reception feedback travels back.
\draw[bgdata] (26.5,66.4)
  .. controls (44,74) and (64,76) .. (85.2,74.6);
\node[bgnote] at (51,75.7) {观测信息回传};
\draw[bgdata] (71.7,63.6)--(86.3,71.5);
\draw[bgdata] (98.5,72.5)--(123,54.5);
\node[bgnote,inner sep=.5mm] at (112,66.5) {信息转发};
\draw[bgfeedback] (139,55)
  .. controls (150,72) and (122,84) .. (98.8,77);
\node[bgnote,inner sep=.5mm] at (130,77.5) {接收确认};

% The ground receiver combines delivered observations for the task decision.
\draw[black!80,fill=white,line width=.7pt]
  (121,44) rectangle (135,54);
\draw[black!40,line width=.4pt] (122.2,46)--(133.8,46);
\draw[black!80,line width=.65pt]
  (128,44)--(128,41.5) (124,41.5)--(132,41.5);
\draw[black!65,line width=.5pt]
  (123,51.8)--(126.5,51.8) (123,49.5)--(130,49.5)
  (131.5,51.8)--(133,51.8);
\draw[black!80,line width=.6pt] (138.5,44)--(138.5,53.3);
\fill[black!80] (138.5,53.7) circle[radius=.6mm];
\draw[black!65,line width=.5pt]
  (139.8,52.5) arc[start angle=-40,end angle=40,radius=1.8mm]
  (140.8,51.7) arc[start angle=-40,end angle=40,radius=3mm];
\node[bgnote] at (128.5,36.8) {地面接收中心};
\node[align=center,font=\fontsize{9bp}{11.5bp}\selectfont]
  at (128.5,26.5) {\textbf{任务目标}\\在时限内汇总信息\\正确判断道路状况};

% Background constraints and their consequences, separate from the methods.
\draw[black!45,line width=.45pt] (1,17.5)--(147,17.5);
\draw[black!25,line width=.4pt] (49.5,6)--(49.5,15.5)
  (98.5,6)--(98.5,15.5);
\node[align=center] at (25,11)
  {\textbf{电量有限、链路变化}\\持续回传受到限制};
\node[align=center] at (74,11)
  {\textbf{视野受限、观测重叠}\\全量上传占用带宽};
\node[align=center] at (123,11)
  {\textbf{数据丢失、反馈延迟}\\按时完成判断更困难};
\node[align=center,font=\fontsize{8bp}{10bp}\selectfont,text=black!75]
  at (74,1.5) {点线：观测范围\qquad 实线箭头：观测信息\qquad 虚线箭头：接收确认};
\end{tikzpicture}
